\section{Aspectos destacados de las diapositivas}

\subsection{Extractos visuales}

Jim Fan enumeró en su guion una serie de diapositivas clave; a continuación se seleccionan páginas representativas para su lectura:

\begin{figure}[p]
\centering
\includegraphics[width=0.95\textwidth,page=2]{slides.pdf}
\caption{Marco semántico y biblioteca de habilidades de Foundation Agent.}
\end{figure}

\begin{figure}[p]
\centering
\includegraphics[width=0.95\textwidth,page=3]{slides.pdf}
\caption{Pipeline de percepción multimodal de GR00T (Vision + Language + Actuation).}
\end{figure}

\begin{figure}[p]
\centering
\includegraphics[width=0.95\textwidth,page=5]{slides.pdf}
\caption{Esquema de domain randomization y guard rail para el Sim-to-Real alignment.}
\end{figure}

\begin{figure}[p]
\centering
\includegraphics[width=0.95\textwidth,page=7]{slides.pdf}
\caption{Operational Playbook y flujo de reproducción de incident trace.}
\end{figure}

\begin{figure}[p]
\centering
\includegraphics[width=0.95\textwidth,page=9]{slides.pdf}
\caption{Ejemplo de Governance checklist y usage policy.}
\end{figure}

\begin{figure}[p]
\centering
\includegraphics[width=0.95\textwidth,page=11]{slides.pdf}
\caption{Multi-robot stack de Project GR00T y su collab plan.}
\end{figure}

\begin{figure}[p]
\centering
\includegraphics[width=0.95\textwidth,page=12]{slides.pdf}
\caption{Ilustración del pipeline que combina LLM + physics solver.}
\end{figure}

\begin{figure}[p]
\centering
\includegraphics[width=0.95\textwidth,page=14]{slides.pdf}
\caption{Replay buffer routing y prioritize schema.}
\end{figure}

\begin{figure}[p]
\centering
\includegraphics[width=0.95\textwidth,page=17]{slides.pdf}
\caption{Trace-driven observability board, combinado con metrics dashboard.}
\end{figure}

\begin{figure}[p]
\centering
\includegraphics[width=0.95\textwidth,page=19]{slides.pdf}
\caption{Detalles del flujo de trabajo Pause + analyze y safety alert icons.}
\end{figure}

\begin{figure}[p]
\centering
\includegraphics[width=0.95\textwidth,page=23]{slides.pdf}
\caption{Future research roadmap: LLM physics, self-supervised, multi-agent.}
\end{figure}

\begin{figure}[p]
\centering
\includegraphics[width=0.95\textwidth,page=25]{slides.pdf}
\caption{Incident response sequence: detect, pause, analyze, recover.}
\end{figure}

\begin{figure}[p]
\centering
\includegraphics[width=0.95\textwidth,page=27]{slides.pdf}
\caption{Reward decomposition chart showing language + physics components.}
\end{figure}

\begin{figure}[p]
\centering
\includegraphics[width=0.95\textwidth,page=30]{slides.pdf}
\caption{Logging taxonomy: sensors, policies, overrides.}
\end{figure}

\begin{figure}[p]
\centering
\includegraphics[width=0.95\textwidth,page=32]{slides.pdf}
\caption{Benchmark ladder comparing dexterity, reliability, generalization.}
\end{figure}

\begin{figure}[p]
\centering
\includegraphics[width=0.95\textwidth,page=34]{slides.pdf}
\caption{Toolchain health dashboard with safety interlocks.}
\end{figure}

\begin{figure}[p]
\centering
\includegraphics[width=0.95\textwidth,page=35]{slides.pdf}
\caption{Data lineage schematic showing robot teleop -> replay -> offline RL.}
\end{figure}

\begin{figure}[p]
\centering
\includegraphics[width=0.95\textwidth,page=36]{slides.pdf}
\caption{Governance maturity curve with audit trail metrics.}
\end{figure}

\begin{figure}[p]
\centering
\includegraphics[width=0.95\textwidth,page=37]{slides.pdf}
\caption{Monitoring dashboard outlining trace, metrics, alerts.}
\end{figure}

\begin{figure}[p]
\centering
\includegraphics[width=0.95\textwidth,page=38]{slides.pdf}
\caption{Incident review template with pause/analyze metrics.}
\end{figure}

\begin{figure}[p]
\centering
\includegraphics[width=0.95\textwidth,page=39]{slides.pdf}
\caption{Lifecycle audit trail showing dataset, config, release pairings.}
\end{figure}

\subsection{Resumen de la sección}

Las diapositivas visualizan la arquitectura de GR00T, la estrategia de simulación, el playbook operativo y el checklist de gobernanza; son material importante para concretar el contenido del texto.

\newpage

